%%=============================================================================
%% PDF-pack Conceptblok B (levels 4-6)
%%
%% Standalone document dat aan deelnemers wordt overhandigd in de "traditionele"
%% conditie van de gebruikersstudie (zie hoofdstuk Evaluatie van de bachelorproef).
%%
%% Compileren: xelatex pdfpack-blokB.tex
%%=============================================================================

\documentclass[11pt,a4paper]{article}

\usepackage[dutch]{babel}
\usepackage[utf8]{inputenc}
\usepackage[T1]{fontenc}
\usepackage[margin=2.5cm]{geometry}
\usepackage{xcolor}
\usepackage{listings}
\usepackage{enumitem}
\usepackage{titlesec}
\usepackage{fancyhdr}
\usepackage{tabularx}
\usepackage{amssymb}
\usepackage{hyperref}

\definecolor{cobolblue}{RGB}{0,80,140}
\definecolor{lightgray}{RGB}{240,240,240}

\lstset{
    basicstyle=\ttfamily\small,
    backgroundcolor=\color{lightgray},
    frame=single,
    framesep=4pt,
    framerule=0pt,
    rulecolor=\color{lightgray},
    breaklines=true,
    columns=fullflexible,
    keywordstyle=\color{cobolblue}\bfseries,
    morekeywords={IDENTIFICATION,DIVISION,PROGRAM-ID,ENVIRONMENT,DATA,WORKING-STORAGE,
        SECTION,PROCEDURE,STOP,RUN,DISPLAY,PIC,PICTURE,VALUE,MOVE,IF,ELSE,END-IF,
        PERFORM,VARYING,UNTIL,END-PERFORM,COMPUTE,FILE,STATUS,SELECT,ASSIGN,FD,
        OPEN,READ,CLOSE,INPUT,OUTPUT,THRU,EVALUATE,WHEN}
}

\pagestyle{fancy}
\fancyhf{}
\fancyhead[L]{\textsc{Learn COBOL \textbar\ PDF-pack Conceptblok B}}
\fancyhead[R]{Pagina \thepage}
\fancyfoot[C]{Bachelorproef Silas De Vuyst, HOGENT, 2025--2026}

\titleformat{\section}{\Large\bfseries\color{cobolblue}}{\thesection}{1em}{}
\titleformat{\subsection}{\large\bfseries}{\thesubsection}{1em}{}

\hypersetup{colorlinks=true,linkcolor=cobolblue,urlcolor=cobolblue}

\begin{document}

%%---------- Titelpagina -----------------------------------------------------
\begin{titlepage}
\begin{center}
\vspace*{2cm}
{\Huge\bfseries Conceptblok B}\\[0.5cm]
{\Large Conditionele logica, iteraties en paragrafen}\\[2cm]

{\large \textit{Traditioneel leermateriaal voor de gebruikersstudie}}\\[0.3cm]
{\large Levels 4 tot en met 6}\\[3cm]

\rule{\textwidth}{0.4pt}\\[0.5cm]
\begin{minipage}{0.9\textwidth}
\textbf{Voor de deelnemer:}
Lees dit pakket aandachtig door en los daarna de opdracht op pagina~\pageref{sec:opdracht-B} op. 
Je hebt 20~minuten om alle leerdoelen op de checklist (\,p.~\pageref{sec:rubric-B})\, te behalen. 
Schrijf je code op papier of in een editor naar keuze, zonder gebruik van automatische validatietools. 
Markeer telkens een nieuwe ``poging'' wanneer je je code aanpast en opnieuw wilt controleren.
\end{minipage}\\
\rule{\textwidth}{0.4pt}\\[1cm]

\vfill
\textbf{Bachelorproef Toegepaste Informatica} \\
Silas De Vuyst \textbullet\ HOGENT \textbullet\ academiejaar 2025--2026
\end{center}
\end{titlepage}

%%---------- Inhoud ----------------------------------------------------------
\section{Inleiding}

In conceptblok~A leerde je hoe een COBOL-programma is opgebouwd uit \textit{divisions} en hoe je data declareert via \texttt{PIC} en level numbers. 
In dit blok ga je een stap verder: je leert COBOL-code \emph{laten beslissen} en \emph{laten herhalen}, en je structureert procedurele logica via benoemde paragraphs.

De drie kernconcepten zijn:

\begin{enumerate}[label=(\arabic*)]
    \item \textbf{Conditionele logica} met \texttt{IF}, \texttt{ELSE} en \texttt{END-IF},
    \item \textbf{Iteraties} via \texttt{PERFORM VARYING \ldots\ UNTIL},
    \item \textbf{Paragrafen} en hun aanroep met een eenvoudige \texttt{PERFORM}.
\end{enumerate}

Voor elke constructie staat naast het COBOL-fragment een Python-equivalent, zodat je het idee aan een vertrouwd patroon kan koppelen.

\section{Level 4: Conditionele logica}

In Python bepaalt indentatie waar een \texttt{if}-blok eindigt. 
COBOL gebruikt expliciete terminators: \texttt{IF \ldots\ ELSE \ldots\ END-IF}.

\textbf{Let op de scope-valstrik.} In oudere COBOL-standaarden beëindigt een \emph{punt} alle openstaande \texttt{IF}-statements. 
Een per ongeluk geplaatste punt kan dus stilletjes je conditionele logica afkappen. 
Werk daarom altijd met \texttt{END-IF} om de scope expliciet af te sluiten.

\subsection*{Voorbeeld}

\begin{lstlisting}
# Python
if age > 18:
    print("Adult")
else:
    print("Minor")
\end{lstlisting}

\begin{lstlisting}
       IDENTIFICATION DIVISION.
       PROGRAM-ID. CONDITIONAL-VB.

       DATA DIVISION.
       WORKING-STORAGE SECTION.
       01 AGE PIC 9(03) VALUE 0.

       PROCEDURE DIVISION.
       MAIN.
           IF AGE > 18
               DISPLAY "Adult"
           ELSE
               DISPLAY "Minor"
           END-IF.
           STOP RUN.
\end{lstlisting}

\section{Level 5: Iteraties}

Waar Python iteratoren gebruikt (\texttt{for i in range(\ldots)}), gebruikt COBOL het \texttt{PERFORM}-statement.
Voor een teller-gestuurde lus is dat \texttt{PERFORM VARYING \ldots\ FROM \ldots\ BY \ldots\ UNTIL \ldots}, afgesloten met \texttt{END-PERFORM}.

\subsection*{Voorbeeld}

\begin{lstlisting}
# Python
for i in range(1, 11):
    print(i)
\end{lstlisting}

\begin{lstlisting}
       IDENTIFICATION DIVISION.
       PROGRAM-ID. LOOP-VB.

       DATA DIVISION.
       WORKING-STORAGE SECTION.
       01 I PIC 9(02) VALUE 1.

       PROCEDURE DIVISION.
       MAIN.
           PERFORM VARYING I FROM 1 BY 1
               UNTIL I > 10
               DISPLAY I
           END-PERFORM.
           STOP RUN.
\end{lstlisting}

\textbf{Belangrijk.} De controlevariabele (hier \texttt{I}) moet altijd vooraf gedeclareerd zijn in de \texttt{DATA DIVISION}.

\section{Level 6: Paragrafen en \texttt{PERFORM}}

Klassieke COBOL kent geen \texttt{def}-functies zoals Python. 
In de plaats werk je met \textbf{paragrafen}: benoemde blokken in de \texttt{PROCEDURE DIVISION}, die je elders aanroept met \texttt{PERFORM <naam>.}

Een paragraafnaam staat alleen op een lijn (kolom~8 of verder), eindigt met een punt en wordt gevolgd door één of meerdere statements.

\subsection*{Voorbeeld}

\begin{lstlisting}
# Python
def add_numbers(a, b):
    return a + b

result = add_numbers(5, 3)
print(result)
\end{lstlisting}

\begin{lstlisting}
       IDENTIFICATION DIVISION.
       PROGRAM-ID. FUNCTION-EX.

       DATA DIVISION.
       WORKING-STORAGE SECTION.
       01 NUM-1  PIC 9(03) VALUE 5.
       01 NUM-2  PIC 9(03) VALUE 3.
       01 RESULT PIC 9(04) VALUE 0.

       PROCEDURE DIVISION.
       MAIN.
           PERFORM ADD-NUMBERS
           DISPLAY RESULT
           STOP RUN.

       ADD-NUMBERS.
           COMPUTE RESULT = NUM-1 + NUM-2.
\end{lstlisting}

%%---------- Opdracht --------------------------------------------------------
\newpage
\section{Opdracht voor conceptblok B}
\label{sec:opdracht-B}

Schrijf één COBOL-programma met de naam \texttt{KLANT-CHECK} dat \emph{alle} onderstaande leerdoelen tegelijkertijd realiseert. 
Je hoeft de code niet uit te voeren: de evaluatie gebeurt op basis van de checklist hieronder.

\begin{enumerate}
    \item Bevat de minimaal vereiste programmastructuur (\texttt{IDENTIFICATION}, \texttt{DATA}, \texttt{PROCEDURE DIVISION}) en eindigt met \texttt{STOP RUN.}
    \item Bevat een conditionele tak met \texttt{IF}, \texttt{ELSE} en expliciete \texttt{END-IF}.
    \item Bevat een lus met \texttt{PERFORM VARYING \ldots\ UNTIL} en expliciete \texttt{END-PERFORM}.
    \item Definieert minstens één benoemde paragraaf naast \texttt{MAIN}.
    \item Roept die paragraaf minstens één keer aan via een \texttt{PERFORM}-statement (anders dan \texttt{PERFORM VARYING}).
\end{enumerate}

\textbf{Voorgesteld scenario.} Laat het programma een klantenlijst doorlopen (lus van 1 t.e.m.\ N), beoordeel per iteratie of een klant ``meerderjarig'' is (\texttt{IF AGE > 18}) en delegeer de afdrukactie naar een paragraaf \texttt{TOON-KLANT}. Andere scenario's zijn ook toegelaten zolang alle leerdoelen aanwezig zijn.

\textbf{Praktisch.} Markeer iedere keer dat je je code controleert tegen de uitleg en aanpast als één ``poging'' en houd zelf bij hoeveel pogingen je nodig had.

%%---------- Rubric ----------------------------------------------------------
\newpage
\section{Zelfcontrole: rubric}
\label{sec:rubric-B}

\begin{tabularx}{\textwidth}{|c|X|}
\hline
\textbf{\#} & \textbf{Leerdoel} \\
\hline
$\square$ & 1. Minimale programmastructuur aanwezig en eindigt met \texttt{STOP RUN.} \\
\hline
$\square$ & 2. \texttt{IF}, \texttt{ELSE} en \texttt{END-IF} correct gebruikt. \\
\hline
$\square$ & 3. \texttt{PERFORM VARYING \ldots\ UNTIL \ldots\ END-PERFORM} correct gebruikt. \\
\hline
$\square$ & 4. Minstens één benoemde paragraaf naast \texttt{MAIN} gedefinieerd. \\
\hline
$\square$ & 5. Minstens één \texttt{PERFORM <paragraaf>} (geen \texttt{PERFORM VARYING}). \\
\hline
\end{tabularx}

\bigskip
\noindent\textbf{Aantal pogingen:}\ \rule{2cm}{0.4pt} \hfill \textbf{Tijd-tot-succes:}\ \rule{2cm}{0.4pt}~min

\bigskip
\noindent\textit{Dank voor je deelname.} Lever dit pakket en je oplossing in bij de onderzoeker.

\end{document}
